% ************************************************************
% INTRODUCTION
% ************************************************************

\section{Introduction}

Conflict-affected WASH assessment requires a manuscript structure that preserves the distinction between household conditions, localized operational observations, infrastructure-system damage, monitoring trends, and health-risk mechanisms. In Gaza, those evidentiary layers are especially vulnerable to overstatement because a single report may describe only a bounded assessment frame, a specific enumeration area, a restricted monitoring period, or an infrastructure recovery estimate rather than a Gaza-wide population parameter.

This study therefore frames Gaza WASH disruption as a source-gated public-health assessment problem. The analytic unit is not an unsupported estimate of total disease burden or infrastructure loss. The analytic unit is a traceable claim connected to a source identifier, role classification, scope limitation, and manuscript-use rule. This design allows household WASH findings, local hotspot evidence, indicative monitoring reports, infrastructure-recovery assessments, and background health-mechanism literature to contribute to the manuscript without collapsing their evidentiary boundaries.

The working research objective is to convert existing Gaza WASH assessment materials into a publication-ready structure in which every manuscript claim is auditable. The central contribution is a reproducible evidence-control layer that can support conflict-health writing while reducing three common errors: treating local factsheets as Gaza-wide evidence, treating infrastructure damage estimates as household outcome surveys, and treating background mechanism papers as current field validation.

% ************************************************************
% END INTRODUCTION
% ************************************************************
